\originalchapter{现代人工智能的问题与证据}\label{ch:route}
\originalsection{从一个预测器到一个工作系统}
一个模型会回答问题, 会用工具, 还会在失败后重新尝试. 这三种行为看起来连成了“自主智慧”, 但要理解近年的进步, 必须先把系统中发生变化的对象分开. 否则, 一次提示变动、一次训练更新、增加一次搜索和增加一道权限检查, 都可能被含混地称为“模型变聪明”. 这些改变需要不同的实验来检验.

本书的学习对象是近年的语言模型与基础模型. 起点不是先学完传统机器学习的全部算法, 而是一个具体问题: 文本序列的下一项预测为什么能够成为回答、推理和使用工具的基础? 阅读顺序沿着这个问题展开. 先看条件概率怎样变成损失和计算图, 再看表示、数据与计算怎样改变性能, 然后比较后训练与推理时计算, 最后把模型放进具有外部动作的研究和工具系统. 矩阵、概率和微分工具附在书后, 在第一次需要时回查相应小节.

\begin{definition}
设 $x$ 是当前上下文, $y$ 是候选输出, $\theta$ 是模型参数. 条件模型 $p_\theta(y\mid x)$ 给输出分配概率. 训练改变 $\theta$; 提示或上下文学习改变 $x$; 采样、搜索和验证改变从该分布得到最终答案的程序; 工具调用改变系统能观察和执行的动作. 这些对象即使联合工作, 也不等同于同一种更新.
\end{definition}
例如, 给固定模型补入几条示范而获得正确回答, 并不证明参数已经学习了新任务. 同样, 把一个错误答案交给精确程序检验后重试, 可能提高最终正确率, 却不能据此说基础模型第一次输出的分布已经改变. 第一次读研究论文时, 先找出作者改变的是参数、上下文、推理程序还是外部系统, 再理解其方法.

工具代理可以写成一个反复运行的闭环. 第 $t$ 步的历史为 $h_t$, 模型提出动作 $a_t\sim p_\theta(\cdot\mid h_t)$, 执行器判断动作是否有权限, 环境返回观察 $o_{t+1}$, 历史更新为 $h_{t+1}=(h_t,a_t,o_{t+1})$. 这个表示允许讨论多步行为, 但没有把循环本身当作意识或意图的证明. 执行器、停止规则、记忆和评价器都是系统的一部分; 只比较模型名称会漏掉它们.

\originalsection{进步需要拆成可比较的问题}
“以前语料差就回答差, 现在为什么还能推理?”包含至少三个问题. 一是预训练数据和目标怎样影响条件分布; 二是同一模型能否把上下文中的任务信息用于未见组合; 三是反馈、搜索和工具怎样改变最后交付的结果. 数据不会失去作用, 但数据的作用不能用“输入好所以输出好”的一句话解释. 数据选择、去重、训练预算和评价污染必须分别控制.

“以前容易被提示注入, 现在为什么更稳?”是另一个问题. 如果模型读到一段网页, 该文本既可能包含任务事实, 也可能包含要求它改写目标的句子. 模型是否能区别来源和权限, 执行器是否允许危险动作, 测试中的攻击者拥有多少尝试机会, 都会影响结果. 指令层级训练与浏览器防御研究直接讨论这些对象\cite{ih,ihblog,browserdefenses}. 较低的攻击成功率提供了指定威胁模型中的鲁棒性证据, 并不自动推出更高的数学能力, 更不推出主观体验.

\begin{definition}
一次能力报告至少指定六个对象: 任务分布 $\mathcal D$, 模型与系统 $\mathcal S$, 可用上下文 $\mathcal C$, 训练及推理预算 $\mathcal B$, 评价规则 $\mathcal E$, 比较基线 $\mathcal S_0$. “系统进步”指在这些条件清楚的比较中, 某个预先指定指标得到改善.
\end{definition}
指标往往形成一个向量, 例如测试负对数似然、答案正确率、攻击成功率、正常任务完成率和计算成本. 一个系统可以在正确率上提高, 在成本或正常任务完成率上下降. 因此不能把若干指标挑选后压成一个没有定义的“智慧分数”. 研究的第一步是把原本混在一起的判断拆成能观察的比较.

\begin{example}
一个固定回答器单次答对的概率为 $p=0.2$. 另一个系统用相同回答器独立采样五次, 并用不出错的验证器挑选正确答案. 最终系统成功率是多少? 这个提升是否证明参数发生了变化?
\end{example}
\begin{solution}
五次都不正确的概率为 $(1-p)^5$, 所以至少一次正确的概率为 $1-(1-p)^5=0.67232$. 这里用到了相同的单次分布、采样独立和验证器无误三个条件. 参数可以完全不变, 改进来自额外计算和选择程序. 如果样本高度相关, 或验证器把错误答案误判为正确, 这个数值不能直接使用. 推理时计算章节将分别处理这些条件.
\end{solution}
这不是某个大模型的实测结果, 而是有限概率模型的算例. 它解释了一种可能的改善机制, 不能替代真实系统的实验. 本书在每次借助小实验解释大系统时都保留这个边界.

\originalsection{阅读顺序与能力递进}
第一阶段围绕条件语言模型. 读者应能写出序列概率的乘积, 解释输入和预测目标为什么错开一位, 追踪一次 token 损失到参数更新, 判断因果掩码和缓存是否正确. 注意力不是从天而降的“智能模块”: 它是可微的内容匹配与信息混合, 需要与位置、残差和训练目标一起理解. 上下文学习首先是固定参数下的条件行为; 哪些内部计算实现了它, 是另一个机制问题.

第二阶段围绕进步证据. 数据处理研究和计算分配研究回答不同问题: 前者比较在其他条件尽量一致时改变语料的效果, 后者比较怎样分配模型、token 和计算预算. 读者应能识别消融、计算是否匹配、测试是否受到污染, 并从一个近似损失模型推导预算分配, 同时说明经验拟合不能成为普遍能力定理\cite{fineweb,chinchilla}.

第三阶段围绕目标与计算的改变. SFT、奖励模型、DPO 与可验证奖励训练改变学习信号; 推理时采样、修订、搜索与工具调用改变运行过程. 阅读 DPO 的意义在于把偏好模型、KL 正则和概率比之间的推导连起来. 阅读 R1 与 RLVR 能力边界研究的意义在于比较单次成功、较大采样预算下的覆盖, 以及训练前后候选分布到底改变了什么. 它们不是简单的“支持 RL”与“反对 RL”两阵营\cite{dpo,r1nature,rlvrboundary,testcompute}.

第四阶段围绕可信行为. 一个代理即使能正确求解题目, 也可能错误执行网页中的动作请求. 相反, 一个执行器可以严格拦截越权动作, 却仍得到错误答案. IH-Challenge 与浏览器防御从来源和动作的信任边界研究鲁棒性; 机制解释则尝试研究内部表示是否被实际使用. 能从表示中解码答案, 不等于该表示因果地控制了输出. CHIVE 提醒读者把“解释看起来合理”转成对新反事实行为的预测测试\cite{chive}.

第五阶段围绕独立研究. 程序搜索与代理自改进需要生成候选、评价、保留档案和盲测. 冻结模型加更好的搜索系统, 与更新基础模型权重, 是两种不同的改进. FunSearch、AlphaEvolve 与 DGM 提供了可以比较的系统实例; 它们的有限任务收益不能直接推出无限自我改进. 新架构也应放进同一个比较: 改变了记忆、计算或生成方式后, 哪个任务改善, 哪个代价下降, 哪项证据仍缺失\cite{funsearch,alphaevolve,dgm}?

\originalsection{怎样读一篇方法论文}
读方法论文时, 先用一页纸记录研究问题、基线和作者实际改变的变量. 然后重建最小计算对象. 对 DPO, 这个对象可以是一组有限动作及其偏好概率; 对注意力, 可以是两位置的因果矩阵; 对提示注入, 可以是带来源标签的观察与一个权限执行器. 缩小对象的目的不是冒充原系统复现, 而是让关键推导和失效条件可以亲手检查.

接着逐一检查三个连接. 第一个连接从方法到目标: 损失到底奖励了什么, 哪些行为只是副作用? 第二个连接从目标到实现: 梯度、掩码、归一化和样本权重是否对应公式? 第三个连接从实现到结论: 指标是否回答研究问题, 基线是否公平, 是否使用测试集选择方法, 结果是否只有一个种子或一个案例? 有了这三个连接, 阅读论文不再等同于记住其摘要.

\begin{definition}
本书区分四类陈述. 数学命题有明确对象和假设, 由证明支持. 作者实验结果是论文在指定设置中的报告. 机制假说解释为什么某种现象发生, 需要有辨别力的检验. 未解决问题是在当前证据下仍不能确定的判断. 本项目的 CPU 数值结果另注明执行程序、输入和环境.
\end{definition}
同一句话可能需要拆开分类. 例如有限动作的 KL 正则最优分布可以严格推导; 用偏好标签训练大模型得到更好回答是指定实验的经验结果; “模型因此形成了新的推理机制”还需要额外证据. 本书的中文译本保留原作者的说法和条件, 译者补充解释另行标注. 原创研究说明则明确列出哪些结论来自原文, 哪些推导和实验由本项目提供.

\originalsection{从复算到研究问题}
研究问题应给出失败判据. 如果希望研究“训练改善的是候选质量还是候选覆盖”, 可以预先固定题集、采样温度和预算, 分别统计单次正确率、pass@$k$、错误类型与验证成本, 再比较训练前后的分布. 如果希望研究“指令层级训练是否泛化到新的攻击”, 应按攻击机制划分测试, 同时衡量正常任务完成率, 并说明攻击者是否自适应. 这些设计都比只追求一个更高总分更接近研究.

本书实验全部可在 CPU 上运行. 它们包含精确可检验的公式与失败案例, 但不声称训练了前沿大模型. 后续研究可以扩大模型、任务和预算, 保留相同的记录原则: 先写问题与比较, 再运行; 保留所有预定条件的结果, 不只保存最好的种子; 对无法复现的结果记录差异, 不把失败删掉后称方法稳定.

\originalsection{练习与解答}
\begin{exercise}
一个代理新增外部检索和验证器之后, 数学题成功率从一个模型基线的报告值提高. 列出至少四项需要核对的条件, 才能判断提升来自哪一种改变.
\end{exercise}
\begin{solution}
应核对模型参数和版本是否相同, 上下文是否增加了答案或额外事实, 总采样/token/工具预算是否匹配, 验证器是否使用了真实答案以及误判率, 题集和划分是否相同, 是否用测试结果选择系统. 参数不变且预算增加时, 首先得到的是系统级计算收益证据. 若同时换了模型、数据和验证器, 不能把收益唯一归于检索.
\end{solution}
\begin{exercise}
设一个 $m$ 步任务的第 $t$ 步在前面各步正确条件下成功的概率为 $q_t$. 不假设各步独立, 写出整条任务成功的概率. 如果只知道每步的无条件成功率, 能否直接相乘?
\end{exercise}
\begin{solution}
令 $A_t$ 表示第 $t$ 步成功. 按条件概率链式法则,
\[
 P\Bigl(\bigcap_{t=1}^{m}A_t\Bigr)
 =P(A_1)\prod_{t=2}^{m}P(A_t\mid A_1\cap\cdots\cap A_{t-1})
 =\prod_{t=1}^{m}q_t.
\]
这是条件概率恒等式, 不需要独立性. 无条件的各步成功率忽略了前一步对后一步输入的影响, 一般不能直接相乘. 多步代理的错误传播与任务级评价因此需要记录完整轨迹.
\end{solution}
\begin{exercise}
“攻击成功率下降”和“模型形成了自主智慧”之间缺少哪些推理? 怎样把前一个结论写成可核查的研究报告?
\end{exercise}
\begin{solution}
后一个说法没有给出智慧的操作定义, 也没有排除训练分布改变、权限执行器和攻击预算等解释. 可核查的写法应给模型/系统版本、攻击面、攻击者权限与预算、任务和攻击采样规则、成功判定、正常任务完成率、比较基线以及不确定性. 这支持所声明条件下的鲁棒性比较, 而不支持关于主观体验的结论.
\end{solution}
\begin{exercise}
一个研究者先试了 100 种评测提示, 只公布测试正确率最高的提示. 即使原模型参数没有改变, 哪个环节已经使用了测试信息? 如何改进设计?
\end{exercise}
\begin{solution}
提示选择已把测试正确率作为反馈, 最终测试不再是独立评价. 可以用开发集选择提示并冻结全部程序, 之后在未使用的测试集评价; 若要研究提示选择本身, 还应报告候选数、预算、所有预定比较和重复试验. 将搜索程序视为系统的一部分, 才能公平讨论它的收益和代价.
\end{solution}
\label{ch:end:01}
